% =====================================================
% 题型：立体几何多选题 (q_solid_prism_properties.tex)
% =====================================================
% =====================================================
% 难度：★★★ (难题)
% =====================================================
\newcommand{\qSolidPrismPropertiesStem}{%
  在直三棱柱 $ABC - A_1B_1C_1$ 中，$AB = 3$，$BC = 4$，$AC = 5$，$AA_1 = 3$，点 $M$ 是棱 $BB_1$ 的中点，下列说法正确的有%
}

\newcommand{\qSolidPrismPropertiesOptions}{%
  \mcoptions[1]{
    $AM \perp BC$
  }{
    四棱锥 $M - ACC_1A_1$ 的体积为 $6$
  }{
    直三棱柱 $ABC - A_1B_1C_1$ 外接球的表面积为 $34\pi$
  }{
    $MA_1 + MC$ 的最小值为 $\sqrt{58}$
  }%
}

\newcommand{\qSolidPrismPropertiesFigure}[1][1.0]{%
  \begin{tikzpicture}[scale=#1, >=stealth]
    % Coordinates for prism
    \coordinate (B) at (0,0);
    \coordinate (C) at (3.0,0);
    \coordinate (A) at (-1.5,1.2);
    
    \coordinate (B1) at (0,3);
    \coordinate (C1) at (3.0,3);
    \coordinate (A1) at (-1.5,4.2);
    
    % Midpoint M
    \coordinate (M) at ($(B)!0.5!(B1)$);
    
    % Visible edges
    \draw[thick] (A) -- (B) -- (C);
    \draw[thick] (A1) -- (B1) -- (C1) -- cycle;
    \draw[thick] (A) -- (A1);
    \draw[thick] (B) -- (B1);
    \draw[thick] (C) -- (C1);
    
    % Hidden edge AC
    \draw[thick, dashed] (A) -- (C);
    
    % Line AM
    \draw[thick] (A) -- (M);
    
    % Points
    \foreach \p in {A,B,C,A1,B1,C1,M} {
      \fill (\p) circle (1.2pt);
    }
    
    % Labels
    \node[left] at (A) {\small $A$};
    \node[below] at (B) {\small $B$};
    \node[right] at (C) {\small $C$};
    \node[left] at (A1) {\small $A_1$};
    \node[above] at (B1) {\small $B_1$};
    \node[right] at (C1) {\small $C_1$};
    \node[right] at (M) {\small $M$};
    
    \ifteacher
      % Draw MC and MA1 for option D visualization
      \draw[thick, gray, dashed] (M) -- (C);
      \draw[thick, gray, dashed] (M) -- (A1);
    \fi
  \end{tikzpicture}%
}

\newcommand{\qSolidPrismPropertiesAnswer}{A, C, D}

\newcommand{\qSolidPrismPropertiesSolution}{%
  对各选项进行分析判定：
  
  对于 \textbf{A} 选项：
  在直三棱柱 $ABC - A_1B_1C_1$ 中，侧棱 $BB_1 \perp$ 底面 $ABC$。
  因为 $BC \subset$ 平面 $ABC$，所以 $BC \perp BB_1$。
  又因为在底面 $\triangle ABC$ 中，$AB = 3$，$BC = 4$，$AC = 5$，满足 $AB^2 + BC^2 = AC^2$，所以 $BC \perp AB$。
  由于 $AB \cap BB_1 = B$，所以 $BC \perp$ 平面 $ABB_1A_1$。
  因为点 $M$ 是棱 $BB_1$ 的中点，所以线段 $AM \subset$ 平面 $ABB_1A_1$，从而 $BC \perp AM$，即 $AM \perp BC$，故 \textbf{A} 正确；
  
  对于 \textbf{B} 选项：
  四棱锥 $M - ACC_1A_1$ 的底面是矩形 $ACC_1A_1$，其面积为：
  \[ S_{ACC_1A_1} = AC \times AA_1 = 5 \times 3 = 15. \]
  四棱锥的高即为点 $M$ 到平面 $ACC_1A_1$ 的距离。
  因为 $BB_1 \parallel$ 平面 $ACC_1A_1$，所以点 $M$ 到平面 $ACC_1A_1$ 的距离等于点 $B$ 到直线 $AC$ 的距离 $h_0$。
  在直角 $\triangle ABC$ 中，由面积关系有 $AC \times h_0 = AB \times BC \implies 5 \times h_0 = 3 \times 4 \implies h_0 = 2.4$。
  所以四棱锥的体积为：
  \[ V_{M-ACC_1A_1} = \frac{1}{3} S_{ACC_1A_1} \times h_0 = \frac{1}{3} \times 15 \times 2.4 = 12. \]
  不为 $6$，故 \textbf{B} 错误；
  
  对于 \textbf{C} 选项：
  直三棱柱的外接球半径 $R$ 满足公式：
  \[ R^2 = r_0^2 + \left(\frac{h_{\text{prism}}}{2}\right)^2, \]
  其中 $r_0$ 为底面直角 $\triangle ABC$ 的外接圆半径，$h_{\text{prism}} = AA_1 = 3$。
  因为 $\triangle ABC$ 是直角三角形，外接圆直径为斜边 $AC = 5$，所以 $r_0 = \dfrac{5}{2}$。
  代入得：
  \[ R^2 = \left(\frac{5}{2}\right)^2 + \left(\frac{3}{2}\right)^2 = \frac{25}{4} + \frac{9}{4} = \frac{34}{4} = \frac{17}{2}. \]
  外接球的表面积为：
  \[ S = 4\pi R^2 = 4\pi \times \frac{17}{2} = 34\pi. \]
  故 \textbf{C} 正确；
  
  对于 \textbf{D} 选项：
  求 $MA_1 + MC$ 的最小值。将侧面 $ABB_1A_1$ 与 $BCC_1B_1$ 沿棱 $BB_1$ 展开在同一个平面内。
  在展开平面中，以 $B$ 为原点，射线 $BC$ 方向为 $x$ 轴正方向，射线 $BA$（展开后）方向为负方向？
  我们画出展开矩形：从左到右依次为 $A - B - C$，其总宽度为 $AB + BC = 3 + 4 = 7$，高度为 $AA_1 = 3$。
  在平面直角坐标系中，设：
  \[ B(0,0),\quad C(4,0),\quad B_1(0,3),\quad C_1(4,3). \]
  侧棱 $AA_1$ 展开后位于左侧，其端点坐标为 $A(-3,0)$，$A_1(-3,3)$。
  中点 $M$ 在 $BB_1$ 上，坐标为 $(0, 1.5)$。
  则 $MA_1 + MC$ 的最小值即为展开后点 $A_1(-3,3)$ 与点 $C(4,0)$ 之间的直线距离。
  \[ A_1C = \sqrt{(4 - (-3))^2 + (0 - 3)^2} = \sqrt{7^2 + (-3)^2} = \sqrt{49 + 9} = \sqrt{58}. \]
  直线段 $A_1C$ 与 $BB_1$（即 $y$ 轴）的交点即为满足最小值的点 $M$ 的位置（由于 $M$ 坐标为 $(0, 1.5)$ 且线段 $A_1C$ 经过 $(0, 1.5)$，确实能取到该最小值）。
  故 $MA_1 + MC$ 的最小值为 $\sqrt{58}$，故 \textbf{D} 正确。
  
  综上所述，正确选项为 \textbf{A, C, D}。
}

\newcommand{\qSolidPrismPropertiesDiagnosis}{%
  用于课堂精准变式训练。考查直三棱柱的线面垂直、体积计算、直棱柱外接球半径公式以及侧面展开最短距离最值问题。
}

\newcommand{\qSolidPrismPropertiesTeachingNotes}{%
  \item \textbf{教学核心}：
    \begin{itemize}
      \item A 选项：先引导学生寻找 $BC$ 垂直的两条相交线；
      \item B 选项：强调棱柱侧棱与平行面距离相等的性质，即点 $M$ 的高就是底面三角形的高 $h_0$；
      \item C 选项：总结直棱柱外接球半径公式 $R^2 = r_0^2 + (h/2)^2$ 是一类高频必拿分题；
      \item D 选项：侧面展开法，注意展开后点 $A_1$ 与 $C$ 分立于 $BB_1$ 的两侧，所以连线就是最小值。
    \end{itemize}
}
